\section{Overview}
\label{sec:overview}
\label{sec:problem}

\xhdr{Problem statement.}
We consider programming a reusable strategy for quasi-static manipulation of rigid bodies from \textit{a single visual human demonstration} $D$ and a language description $l$. The program should solve a family of tasks $\mathcal{E}$ that share similar objectives but may vary in object appearance, geometry, pose, physical properties, and scene configuration. No manipulation primitives, explicit success or reward signals, or instantiated simulation environments are provided \textit{a priori}. Given only $(D,l)$, the agent aims to construct executable primitives $\mathcal{P}=\{P_1,\ldots,P_N\}$ and compose them into a strategy $\Pi$ that captures the underlying structure of the demonstration and generalizes to unseen language instructions $l^\star$ and test scenes $E^\star$ in $\mathcal{E}$:
\[
(D,l)
\mapsto
(\mathcal{P},\Pi),
\quad
\Pi(E^\star, l^\star;\mathcal{P}) \text{ succeeds for } (E^\star, l^\star) \in \mathcal{E}.
\]

\xhdr{Overview of \ours.}
As shown in \fig{fig:method_overview}, RAPID closes the agentic coding loop by deriving task specifications, manipulation primitives, and verification environments from $(D,l)$. It partitions the demonstration into interaction phases, infers a subgoal and success predicate for each phase, and reconstructs the corresponding interactive simulation environments. To improve generalization, RAPID also synthesizes feasible scene variants for additional verification. For each phase, the coding agent generates a primitive to achieve its subgoal and iteratively executes, verifies, and revises it in the corresponding environment and its variants. It then infers the overall task goal and success predicate, composes the verified primitives into a strategy, and refines the composition through the same loop. At deployment, given an RGB-D observation of a novel scene and a language instruction, RAPID reconstructs the scene in simulation, instantiates and runs the strategy program in this scene, and executes the output trajectory on the robot. We describe program construction and verification in Sec.~\ref{sec:agentic_programming}.

We introduce a \textit{relational program representation} to make manipulation strategies reusable across scenes. A strategy connects primitives through relational constraints that specify what each primitive should accomplish in the current scene. Each primitive then uses local trajectory optimization to achieve the requested object-level effect, taking the scene's geometry and physical properties into account. To deploy the program in a new scene, we use lightweight semantic binding to match scene objects to program roles. The strategy and primitives themselves remain unchanged. We provide the details of the representation in Sec.~\ref{sec:relational_programs}.
